\documentclass[tikz,border=5pt]{standalone}
\usetikzlibrary{calc,arrows.meta}
\begin{document}
\def\Angle{-28}
\def\Curvature{0.13}
\def\Radius{2.05}
\def\Mesh{0.30}
\def\Eta{0.15} % Boundary-layer half-width is 2\eta
\def\Px{0}
\def\Py{1.65}
\pgfmathsetmacro{\ct}{cos(\Angle)}
\pgfmathsetmacro{\st}{sin(\Angle)}
\pgfmathsetmacro{\umin}{-2.9}
\pgfmathsetmacro{\umax}{2.9}
\pgfmathsetmacro{\rr}{\Radius-0.18}
% Endpoint of the curve inside the ball, computed numerically by bisection.
\pgfmathsetmacro{\ub}{\Radius}
\foreach \i in {1,...,25}{\pgfmathsetmacro{\ub}{0.5*(\ub+\Radius/sqrt(1+0.25*\Curvature*\Curvature*\ub*\ub))}}
\pgfmathsetmacro{\leftX}{\Px-\ub*\ct-0.5*\Curvature*\ub*\ub*\st}
\pgfmathsetmacro{\rightX}{\Px+\ub*\ct-0.5*\Curvature*\ub*\ub*\st}
\pgfmathsetmacro{\ellX}{\Px-\Radius*\ct}
\pgfmathsetmacro{\rX}{\Px+\Radius*\ct}
\pgfmathsetmacro{\rminus}{\rr*\ct}
\pgfmathsetmacro{\fleft}{\leftX-0.11}
\pgfmathsetmacro{\fright}{\rightX+0.11}
\pgfmathsetmacro{\band}{2*\Eta}
\tikzset{curve/.style={black,very thick}, tangent/.style={red!80!black,dashed,thick}, ball/.style={blue!75!black,dashed,thick}, guide/.style={gray!65,densely dotted}, axis/.style={-Latex,thin}}
\begin{tikzpicture}[scale=1]
\foreach \shift/\panel in {0/{(a) Local projections},8.2/{(b) Homotopy}}{
\begin{scope}[xshift=\shift cm]
  \draw[axis] (-3.1,0)--(3.25,0) node[right] {$x$};
  \draw[axis] (-3,-0.2)--(-3,4.25) node[above] {$y$};
  \node[anchor=west,font=\bfseries] at (-3,4.75) {\panel};
  \draw[ball] (\Px,\Py) circle[radius=\Radius];
  \node[blue!70!black] at (1.65,3.55) {$B(p,r)$};
  \draw[tangent] ({\Px+\umin*\ct},{\Py+\umin*\st})--({\Px+\umax*\ct},{\Py+\umax*\st});
  \draw[curve,domain=\umin:\umax,samples=100,smooth,variable=\u] plot ({\Px+\u*\ct-0.5*\Curvature*\u*\u*\st},{\Py+\u*\st+0.5*\Curvature*\u*\u*\ct});
  \fill (\Px,\Py) circle (1.8pt) node[above right] {$p$};
  \node[red!75!black] at (2.45,0.72) {$T_pX$};
\end{scope}
}
% Voxel squares: a grid-based approximation, clipped visually by the ball.
\begin{scope}
\clip (\Px,\Py) circle[radius=\Radius];
\foreach \i in {-10,...,10}{
 \foreach \j in {-4,...,16}{
  \pgfmathsetmacro{\xx}{(\i+0.5)*\Mesh}
  \pgfmathsetmacro{\yy}{(\j+0.5)*\Mesh}
  \pgfmathsetmacro{\uu}{(\xx-\Px)*\ct+(\yy-\Py)*\st}
  \pgfmathsetmacro{\vv}{-(\xx-\Px)*\st+(\yy-\Py)*\ct}
  \pgfmathsetmacro{\distance}{abs(\vv-0.5*\Curvature*\uu*\uu)}
  \ifdim\distance pt<0.22pt
    \filldraw[fill=blue!25,draw=blue!40,line width=.18pt] (\i*\Mesh,\j*\Mesh) rectangle ++(\Mesh,\Mesh);
  \fi
 }
}
\end{scope}
% Redraw curve and tangent for clarity
\draw[tangent] ({\Px+\umin*\ct},{\Py+\umin*\st})--({\Px+\umax*\ct},{\Py+\umax*\st});
\draw[curve,domain=\umin:\umax,samples=100,variable=\u] plot ({\Px+\u*\ct-0.5*\Curvature*\u*\u*\st},{\Py+\u*\st+0.5*\Curvature*\u*\u*\ct});
\fill (\Px,\Py) circle (1.8pt);
% Projected intervals: E tangent disk, F curve projection, blue enlarged voxel envelope.
\draw[guide] (\ellX,\Py-\Radius*\st/2)--(\ellX,-.65);
\draw[guide] (\rX,\Py+\Radius*\st/2)--(\rX,-.65);
\draw[red!75!black,line width=3pt,opacity=.65] (\ellX,-.35)--(\rX,-.35);
\node[red!75!black,below] at (0,-.40) {$\mathcal E_J$};
\draw[blue!75!black,line width=3pt,opacity=.65] (\fleft,-.95)--(\fright,-.95);
\node[blue!75!black,below] at (0,-1.0) {$\mathcal F_J$};
% Boundary of E_J thickened by 2 eta: two neighborhoods of its endpoints.
\foreach \edge in {\ellX,\rX}{
  \fill[green!45,opacity=.48] ({\edge-\band},-1.75) rectangle ({\edge+\band},-1.52);
  \draw[green!45!black,thin] ({\edge-\band},-1.75) rectangle ({\edge+\band},-1.52);
}
\node[green!35!black,anchor=north,font=\scriptsize] at (0,-1.79)
 {$ (\partial\mathcal E_J)_{2\eta}$};
% Symmetric difference: pieces of E_J not in F_J and F_J not in E_J.
% These are calculated from the same endpoints used for the intervals above.
\pgfmathsetmacro{\missingLeft}{max(0,\fleft-\ellX)}
\pgfmathsetmacro{\extraLeft}{max(0,\ellX-\fleft)}
\pgfmathsetmacro{\missingRight}{max(0,\rX-\fright)}
\pgfmathsetmacro{\extraRight}{max(0,\fright-\rX)}
\ifdim\missingLeft pt>0pt
 \fill[orange!85] (\ellX,-1.32) rectangle (\fleft,-1.16);
\fi
\ifdim\extraLeft pt>0pt
 \fill[orange!85] (\fleft,-1.32) rectangle (\ellX,-1.16);
\fi
\ifdim\missingRight pt>0pt
 \fill[orange!85] (\fright,-1.32) rectangle (\rX,-1.16);
\fi
\ifdim\extraRight pt>0pt
 \fill[orange!85] (\rX,-1.32) rectangle (\fright,-1.16);
\fi
\node[orange!75!black,anchor=north,font=\scriptsize] at (0,-1.34)
 {$\mathcal E_J\triangle\mathcal F_J$};
\node[anchor=west,font=\scriptsize] at (-2.9,-2.15)
 {$\mathcal E_J\triangle\mathcal F_J\subset(\partial\mathcal E_J)_{2\eta}$};
% Second panel: homotopy curves use SAME angle and curvature, all tangent at p.
\begin{scope}[xshift=8.2cm]
\foreach \t in {0.25,0.5,0.75}{
\draw[gray!70,line width=.9pt,domain=-\rr:\rr,samples=70,variable=\u] plot ({\Px+\u*\ct-\t*0.5*\Curvature*\u*\u*\st},{\Py+\u*\st+\t*0.5*\Curvature*\u*\u*\ct});
}
\draw[very thick,orange!85!black,domain=-\rr:\rr,samples=70,variable=\u] plot ({\Px+\u*\ct-0.5*\Curvature*\u*\u*\st},{\Py+\u*\st+0.5*\Curvature*\u*\u*\ct});
\draw[very thick,red!70!black] ({\Px-\rr*\ct},{\Py-\rr*\st})--({\Px+\rr*\ct},{\Py+\rr*\st});
\node[anchor=west,fill=white,inner sep=1pt,font=\small] at (-2.9,-.55) {$G_{J,t}(u)=p_x+u\cos\theta-tf(u)\sin\theta$};
\node[gray!70!black] at (2.0,3.1) {$0<t<1$};
\node[orange!80!black] at (2.1,2.8) {$t=1$};
\node[red!75!black] at (2.1,2.5) {$t=0$};
\end{scope}
\end{tikzpicture}
\end{document}
